\section{\texorpdfstring{The Kummer Field and Its Quadratic $L$-Functions}{The Kummer Field and Its Quadratic L-Functions}}\label{gf:section}

The analytic estimate of Sections~\ref{an:section}--\ref{ce:section}
compares every field of Theorem~\ref{tw:field-family} with a fixed subfield.
The simplest choice of this subfield is the Kummer field
\[
 E=B(\sqrt{\alpha_0},\ldots,\sqrt{\alpha_7})=B(\sqrt V)
\]
of Subsection~\ref{tw:global-section}, with $\alpha_0,\ldots,\alpha_7$ as in
\eqref{tw:alpha} and $V$ as in Lemma~\ref{tw:base}(d). It is the maximal
elementary abelian $2$-extension of $B$ unramified at the finite primes
outside $S$, that is, the fixed field of the Frattini subgroup of $G_B$
(Lemma~\ref{gf:in-tower}). The Lean formalization calls $E$ the genus field;
we do not use that name here. Here the comparison field is a larger field
$\EW\supset E$ (Section~\ref{dh:section}), whose zeta function contains
$\zeta_E$ as a factor. This section shows that $E$ lies in every field of
Theorem~\ref{tw:field-family}. It then writes $\zeta_E$ as the product of
$\zeta_B$ and $255$ quadratic Hecke $L$-functions of $B$, and determines
their conductors, gamma factors, root numbers and Euler factors. The field
$E$ does not depend on the choice of the capped primes.

For $v,e\in\Ftwo^8$ put $\langle v,e\rangle=\sum_{i=0}^7v_ie_i\in\Ftwo$; as
in Section~\ref{tw:tower}, vectors are written as strings of their
coordinates. For a nonzero ideal $\mathfrak a$ of a number field,
$\mathrm N\mathfrak a$ denotes its absolute norm. The number of positive
divisors of an integer $n\ge1$ is $\tau(n)$; it is not to be confused with
the inertia generators $\tau_{\mathfrak p}$ of
Definition~\ref{tw:local-generators}, whose elementary images
$\bar\tau_{\mathfrak p}$ appear in Lemma~\ref{gf:linear}.

\subsection{The Characters of the Kummer Field}\label{gf:characters-section}
This subsection fixes the characters $\chi_e$ of $\Gal(E/B)\cong\Ftwo^8$
(Definition~\ref{gf:characters}) and shows that $E$ lies in every field of
the tower (Lemma~\ref{gf:in-tower}).

By Lemma~\ref{tw:kummer-field}, $[E:B]=256$, and the elementary image
$\bar g\in\Ftwo^8$ of Definition~\ref{tw:elementary-image}, given by
$g(\sqrt{\alpha_i})=(-1)^{\bar g_i}\sqrt{\alpha_i}$, identifies $\Gal(E/B)$
with $\Ftwo^8$.

\begin{definition}\label{gf:characters}
For $e\in\Ftwo^8$ put
\[
 \alpha_e=\prod_{i=0}^7\alpha_i^{e_i},\qquad B_e=B(\sqrt{\alpha_e}),
\]
and let $\chi_e$ be the character $g\mapsto(-1)^{\langle\bar g,e\rangle}$ of
$\Gal(E/B)$.
\end{definition}

The character $\chi_e$ takes the values $\pm1$; in terms of the
$\Ftwo$-valued Kummer characters $\kappa_a$ of
Definition~\ref{tw:elementary-image}, $\chi_e(g)=(-1)^{\kappa_{\alpha_e}(g)}$.
Since $g(\sqrt{\alpha_e})=\chi_e(g)\sqrt{\alpha_e}$,
the character $\chi_e$ cuts out $B_e$, and the $256$ characters $\chi_e$ are
all the characters of $\Gal(E/B)$. For $e\ne0$ the extension $B_e/B$ is
quadratic, because the classes of $\alpha_0,\ldots,\alpha_7$ are independent
in $B^\times/B^{\times2}$ (Lemma~\ref{tw:base}(d)).

\begin{lemma}\label{gf:in-tower}
The fixed field of the Frattini subgroup $D_2G_B$ of $G_B$ is $E$.
Consequently $E$ is contained in every field $K$ of
Theorem~\ref{tw:field-family}.
\end{lemma}

\begin{proof}
Let $G_S$ be the group of Subsection~\ref{tw:global-section}, so
that $G_S/D_2G_S=\Gal(E/B)$ (Lemma~\ref{tw:kummer-field}), and let
$G_B=G_S/N$ as in Definition~\ref{tw:GB}. By Lemma~\ref{tw:GB-basic}(b),
$N\subseteq D_2G_S$, so
$G_B/D_2G_B=G_S/D_2G_S=\Gal(E/B)$. A field $K$ of
Theorem~\ref{tw:field-family} contains the fixed field of
$D_4G_B\subseteq D_2G_B$, and hence contains $E$.
\end{proof}

\subsection{Local Data}\label{gf:local-section}
This subsection computes the conductor exponents and the Frobenius values of
every character $\chi_e$ (Proposition~\ref{gf:local-data}), expresses them
through the elementary images of the local generators
(Lemma~\ref{gf:linear}), and records the types of the primes of $B$ in $E/B$
(Corollary~\ref{gf:types}). They give the Euler factors and conductors of
Subsection~\ref{gf:zeta-section}, and the Frobenius vectors used by the
census of Section~\ref{an:section}.

\begin{definition}\label{gf:frobenius-vector}
For a prime $\mathfrak p\notin S$ of $B$, the \emph{Frobenius vector}
$\overline{\Frob}_{\mathfrak p}\in\Ftwo^8$ is the elementary image of the
Frobenius of $\mathfrak p$ in $\Gal(E/B)$, that is, of the Frobenius element
$\Frob_{\mathfrak p}$ of Subsection~\ref{tw:global-section}.
\end{definition}

For $\mathfrak p=\mathfrak t_k$, $k=1,2,3$, it is the vector of
$\Frob_{\mathfrak t_k}$ in Table~\ref{tw:vector-table}, and for
$\mathfrak p=\mathfrak t_4$ it is $00110100$ (Lemma~\ref{tw:vectors}). By
Euler's criterion,
$(\overline{\Frob}_{\mathfrak p})_i=1$ exactly when $\alpha_i$ is not a
square in $\mathcal O_B/\mathfrak p$. For a
place $\mathfrak p$ of $B$ we write $(\cdot,\cdot)_{\mathfrak p}$ for the
quadratic Hilbert symbol of $B_{\mathfrak p}$.

\begin{proposition}\label{gf:local-data}
Let $e\in\Ftwo^8\setminus\{0\}$. For a prime $\mathfrak p$ of $B$ let
$c_{\mathfrak p}(e)$ be the exponent of $\mathfrak p$ in the conductor of
$\chi_e$, and if $c_{\mathfrak p}(e)=0$, let $\chi_e(\mathfrak p)\in\{\pm1\}$
be the value of $\chi_e$ at the Frobenius of $\mathfrak p$.
\begin{enumerate}
\item If $\mathfrak p\notin S$, then $c_{\mathfrak p}(e)=0$ and
 $\chi_e(\mathfrak p)=(-1)^{\langle\overline{\Frob}_{\mathfrak p},e\rangle}$.
 Let $p$ be the rational prime below $\mathfrak p$. If $p$ splits in $B$, then
 $\mathfrak p=(p,\sqrt{241}-r)$ with $r^2\equiv241\pmod p$, and
 $(\overline{\Frob}_{\mathfrak p})_i=1$ exactly when the image of $\alpha_i$
 under $\sqrt{241}\mapsto r$ is a quadratic nonresidue modulo $p$. If $p$ is
 inert, then $(\overline{\Frob}_{\mathfrak p})_i=1$ exactly when
 $\mathrm N_{B/\Q}(\alpha_i)$ is a quadratic nonresidue modulo $p$. If $p=241$, the
 same holds with $\sqrt{241}\mapsto0$.
\item Let $\mathfrak p$ lie above $p\in\{3,5\}$, let $u=-1$ if $p=3$ and $u=2$
 if $p=5$, and let $\alpha_i$ be the generator of $\mathfrak p$. If
 $(u,\alpha_e)_{\mathfrak p}=-1$, then $c_{\mathfrak p}(e)=1$. Otherwise
 $c_{\mathfrak p}(e)=0$ and
 $\chi_e(\mathfrak p)=(-\alpha_i,\alpha_e)_{\mathfrak p}$. The Frobenius
 generator $\varphi_{\mathfrak p}$ of Definition~\ref{tw:local-generators},
 which fixes $\sqrt{\alpha_i}$, acts on square roots by
 $(-\alpha_i,\cdot)_{\mathfrak p}$ (Table~\ref{tw:vector-table}).
\item Let $\mathfrak p=\mathfrak p_j$. If $(5,\alpha_e)_{\mathfrak p}=-1$, then
 $c_{\mathfrak p}(e)=3$. If $(5,\alpha_e)_{\mathfrak p}=1$ and
 $(-1,\alpha_e)_{\mathfrak p}=-1$, then $c_{\mathfrak p}(e)=2$. Otherwise
 $c_{\mathfrak p}(e)=0$ and $\chi_e(\mathfrak p)=(-2,\alpha_e)_{\mathfrak p}$.
\item The real place $v_k$ becomes complex in $B_e$, that is, $\chi_e$ is
 nontrivial at $v_k$, exactly when $\alpha_e<0$ at $v_k$.
\item The conductor of $\chi_e$ is the ideal
 $\mathfrak f_e=\prod_{\mathfrak p\in S}\mathfrak p^{c_{\mathfrak p}(e)}$ of
 $\mathcal O_B$.
\end{enumerate}
\end{proposition}

\begin{proof}
For (1), $\mathfrak p$ is odd and $\alpha_e$ is a unit at $\mathfrak p$, so
$\mathfrak p$ is unramified in $B_e$. By Euler's criterion the Frobenius
multiplies $\sqrt{\alpha_e}$ by the quadratic residue symbol of $\alpha_e$
modulo $\mathfrak p$, which is $\prod_i(\alpha_i/\mathfrak p)^{e_i}$. For a
split prime, $\mathcal O_B/\mathfrak p\cong\mathbb F_p$ via $\sqrt{241}\mapsto
r$. For an inert prime, $\mathcal O_B/\mathfrak p\cong\mathbb F_{p^2}$ and
the conjugation of $B$ induces the Frobenius $x\mapsto x^p$ of the residue
field, so the residue of $\mathrm N_{B/\Q}(\alpha)$ is $\bar\alpha^{p+1}$. An element
of $\mathbb F_{p^2}^\times$ is a square exactly when
$\bar\alpha^{(p^2-1)/2}=(\bar\alpha^{p+1})^{(p-1)/2}$ equals $1$, that is,
when its norm is a square in $\mathbb F_p$. The prime above $241$ has
residue field $\mathbb F_{241}$ and $\sqrt{241}\mapsto0$.

For (2) and (3) we use local class field theory for $B_{\mathfrak p}=\Q_p$.
The local Artin map sends $b\in\Q_p^\times$ to the automorphism multiplying
$\sqrt{\alpha_e}$ by $(b,\alpha_e)_{\mathfrak p}$. Units map onto inertia, a
uniformizer maps to a Frobenius, and the conductor exponent is the least
$n\ge0$ such that $b\mapsto(b,\alpha_e)_{\mathfrak p}$ is trivial on
$U^{(n)}$, where $U^{(0)}=\mathbb Z_p^\times$ and $U^{(n)}=1+p^n\mathbb Z_p$
\cite{MilneCFT}.

For $p=3,5$, the unit group modulo squares is $\{1,u\}$, and $U^{(1)}$
consists of squares by Hensel's lemma. Hence $c_{\mathfrak p}(e)\le1$, with
equality exactly when $(u,\alpha_e)_{\mathfrak p}=-1$. In the unramified
case the Frobenius is the image of any uniformizer, such as $-\alpha_i$.

For $p=2$, the squares of $\mathbb Z_2^\times$ form $U^{(3)}$. Moreover
$U^{(2)}=U^{(3)}\cup5U^{(3)}$ and
$U^{(0)}=U^{(1)}=\{\pm1,\pm5\}U^{(3)}$. The character is thus trivial on
$U^{(3)}$. It is trivial on $U^{(2)}$ exactly when $(5,\alpha_e)=1$, and
on $U^{(0)}$ exactly when also $(-1,\alpha_e)=1$. This gives the stated
exponents; the exponent $1$ cannot occur since $U^{(0)}=U^{(1)}$. In the
unramified case the Frobenius is the image of the uniformizer $-2$.

For (4), the real place $v_k$ splits in $B_e$ exactly when $\alpha_e>0$ at
$v_k$. For (5), $c_{\mathfrak p}(e)=0$ for $\mathfrak p\notin S$ by (1).
\end{proof}

Each condition in the proposition is linear in $e$:

\begin{lemma}\label{gf:linear}
Let $e\in\Ftwo^8$. Then
$(5,\alpha_e)_{\mathfrak p_j}=(-1)^{\langle\bar x_j,e\rangle}$,
$(-1,\alpha_e)_{\mathfrak p_j}=(-1)^{\langle\bar y_j,e\rangle}$ and
$(-2,\alpha_e)_{\mathfrak p_j}=(-1)^{\langle\bar z_j,e\rangle}$. For a prime
$\mathfrak p$ above $3$ or $5$, with $u$ as in
Proposition~\ref{gf:local-data}(2) and $\alpha_i$ the generator of
$\mathfrak p$,
$(u,\alpha_e)_{\mathfrak p}=(-1)^{\langle\bar\tau_{\mathfrak p},e\rangle}$ and
$(-\alpha_i,\alpha_e)_{\mathfrak p}=(-1)^{\langle\bar\varphi_{\mathfrak p},e\rangle}$.
Finally, $\alpha_e<0$ at $v_k$ exactly when $\langle\bar\iota_k,e\rangle=1$.
\end{lemma}

\begin{proof}
In each case $(b,\alpha_e)_{\mathfrak p}=(-1)^{\langle w,e\rangle}$, where $w$
is the elementary image of the local generator that acts on square roots by
$(b,\cdot)_{\mathfrak p}$ in Table~\ref{tw:vector-table}: at $\mathfrak p_j$
the vectors $\bar x_j,\bar y_j,\bar z_j$ belong to $b=5,-1,-2$, and above
$3$ and $5$ the images of the inertia and Frobenius generators belong to
$b=u$ and $b=-\alpha_i$. Likewise the complex conjugation $\iota_k$ negates
$\sqrt{\alpha_e}$ exactly when $\alpha_e<0$ at $v_k$
(Lemma~\ref{tw:local-groups}(a)), and it multiplies $\sqrt{\alpha_e}$ by
$(-1)^{\langle\bar\iota_k,e\rangle}$.
\end{proof}

\begin{definition}\label{gf:archimedean}
For $e\in\Ftwo^8$ and $k=1,2$ put
$\nu_k(e)=\langle\bar\iota_k,e\rangle\in\{0,1\}$, where $\bar\iota_k$ is the
elementary image of the complex conjugation $\iota_k$
(Table~\ref{tw:vector-table}).
\end{definition}

By Proposition~\ref{gf:local-data}(4) and Lemma~\ref{gf:linear}, for $e\ne0$
we have $\nu_k(e)=1$ exactly when the real place $v_k$ becomes complex in
$B_e$, that is, when $\chi_e$ is nontrivial at $v_k$.

\begin{corollary}\label{gf:types}
The type (Definition~\ref{tw:type}) of a prime $\mathfrak p$ of $B$ in
$E/B$ is $(4,2)$ at $\mathfrak p_1,\mathfrak p_2$; $(2,2)$ at
$\mathfrak q_1,\mathfrak q_2,\mathfrak r_1,\mathfrak r_2$; and $(1,1)$ or
$(1,2)$ at $\mathfrak p\notin S$, according as $\overline{\Frob}_{\mathfrak p}=0$
or not.
In particular the capped primes $\mathfrak t_1,\mathfrak t_2,\mathfrak t_3$
have type $(1,2)$ in $E/B$.
\end{corollary}

\begin{proof}
For $\mathfrak p\in S$, local class field theory identifies the
decomposition group of $\mathfrak p$ in $\Gal(E/B)=\Ftwo^8$ with the image
of $\Q_p^\times$ under the local Artin map, and the inertia group with the
image of the units. By Lemma~\ref{gf:linear}, at $\mathfrak p_j$ these are
spanned by $\bar x_j,\bar y_j,\bar z_j$ and by $\bar x_j,\bar y_j$; the three
vectors are independent (Table~\ref{tw:vector-table}), so the decomposition
group has order $8$ and inertia has order $4$. Above $3$ and $5$ the two
elementary images are independent, which gives $(2,2)$. For
$\mathfrak p\notin S$ the decomposition group is generated by
$\overline{\Frob}_{\mathfrak p}$. The Frobenius vectors of $\mathfrak t_1,\mathfrak
t_2,\mathfrak t_3$ are nonzero (Table~\ref{tw:vector-table}).
\end{proof}

\subsection{The Zeta Function of the Kummer Field}\label{gf:zeta-section}
This subsection factors $\zeta_E$ into $\zeta_B$ and $255$ quadratic Hecke
$L$-functions (Proposition~\ref{gf:factorization}), and records their
coefficients, completed forms, functional equations and conductors
(Proposition~\ref{gf:degree-two}, Corollary~\ref{gf:conductors}). These are
the input of the approximate functional equation in
Section~\ref{an:section}.

Let $\chi_B$ be the Dirichlet character $n\mapsto(n/241)$.

\begin{proposition}\label{gf:factorization}
For $\operatorname{Re}s>1$,
\[
 \zeta_E(s)=\zeta_B(s)\prod_{e\in\Ftwo^8\setminus\{0\}}L(s,\chi_e),
 \qquad \zeta_B(s)=\zeta(s)L(s,\chi_B),
\]
where $L(s,\chi_e)=\zeta_{B_e}(s)/\zeta_B(s)$ is the Hecke $L$-function of
$\chi_e$.
\end{proposition}

\begin{proof}
Compare Euler factors at a prime $\mathfrak p$ of $B$. Let $T_{\mathfrak
p}\subseteq Z_{\mathfrak p}\subseteq G=\Gal(E/B)$ be its inertia and
decomposition groups, and put $f_{\mathfrak p}=|Z_{\mathfrak p}/T_{\mathfrak
p}|$ and $g_{\mathfrak p}=|G/Z_{\mathfrak p}|$. The Euler factor of $\zeta_E$
at $\mathfrak p$ is $(1-X^{f_{\mathfrak p}})^{-g_{\mathfrak p}}$ with
$X=\mathrm N\mathfrak p^{-s}$. A character nontrivial on $T_{\mathfrak p}$
contributes the factor $1$. The characters trivial on $T_{\mathfrak p}$ are
those of $G/T_{\mathfrak p}$; they restrict onto the $f_{\mathfrak p}$
characters of the cyclic group $Z_{\mathfrak p}/T_{\mathfrak p}$, each with
$g_{\mathfrak p}$ extensions. Hence their factors multiply to
$\prod_{\zeta^{f_{\mathfrak p}}=1}(1-\zeta X)^{-g_{\mathfrak p}}
=(1-X^{f_{\mathfrak p}})^{-g_{\mathfrak p}}$. The same identity for $B_e/B$
identifies $L(s,\chi_e)$ with $\zeta_{B_e}/\zeta_B$. The second identity is
the case of $B/\Q$. Since $241\equiv1\pmod8$, the prime $2$ splits in $B$,
and quadratic reciprocity shows that the character of $B/\Q$ is
$\chi_B$.
\end{proof}

Put $\Gamma_\R(s)=\pi^{-s/2}\Gamma(s/2)$ and
$\Gamma_\C(s)=2(2\pi)^{-s}\Gamma(s)$, so that
$\Gamma_\C(s)=\Gamma_\R(s)\Gamma_\R(s+1)$ by the duplication formula. For a
number field $M$ with discriminant $\Delta_M$ and signature
$(r_1(M),r_2(M))$, the completed Dedekind zeta function is
\[
 \Lambda_M(s)=|\Delta_M|^{s/2}\Gamma_\R(s)^{r_1(M)}\Gamma_\C(s)^{r_2(M)}
 \zeta_M(s).
\]

\begin{proposition}\label{gf:degree-two}
Let $e\ne0$, write $L(s,\chi_e)=\sum_{n\ge1}a_n(e)n^{-s}$, and put
\[
 Q_e=241\,\mathrm N\mathfrak f_e,\qquad
 \Lambda(s,\chi_e)=Q_e^{s/2}\,\Gamma_\R(s+\nu_1(e))\,\Gamma_\R(s+\nu_2(e))\,
 L(s,\chi_e).
\]
\begin{enumerate}
\item The coefficients $a_n(e)$ are integers with $|a_n(e)|\le\tau(n)$. The
 Euler factor at a rational prime $p$ is
 $\prod_{\mathfrak p\mid p}(1-\chi_e(\mathfrak p)\mathrm N\mathfrak p^{-s})^{-1}$,
 with $\chi_e(\mathfrak p)=0$ when $c_{\mathfrak p}(e)>0$.
\item The function $\Lambda(s,\chi_e)$ equals $\Lambda_{B_e}(s)/\Lambda_B(s)$.
 It is entire of order at most one and satisfies
 $\Lambda(s,\chi_e)=\Lambda(1-s,\chi_e)$; that is, the root number of
 $L(s,\chi_e)$ is $1$.
\item Exactly $63$ of the characters have $\nu_1(e)=\nu_2(e)=0$, exactly
 $64$ have $\nu_1(e)=\nu_2(e)=1$, and $128$ have $\nu_1(e)\ne\nu_2(e)$.
\end{enumerate}
\end{proposition}

\begin{proof}
The Euler factor in (1) is Proposition~\ref{gf:factorization} applied prime
by prime. Above a rational prime there are at most two primes of $B$. So
the local factor is either a product of at most two factors
$(1-cp^{-s})^{-1}$, or a single factor $(1-cp^{-2s})^{-1}$, with
$c\in\{0,\pm1\}$. In both cases the coefficient of $p^{-ks}$ has modulus at
most $k+1=\tau(p^k)$. Multiplicativity gives (1).

For (2), the completed Dedekind zeta functions satisfy
$\Lambda_M(1-s)=\Lambda_M(s)$, and $L(s,\chi_e)$ extends to an entire
function because $\chi_e$ is nontrivial \cite{Hecke1920,Tate1967}. The
discriminants satisfy
$|\Delta_{B_e}|=|\Delta_B|^2\,\mathrm N\mathfrak d_{B_e/B}$, where
$\mathfrak d_{B_e/B}$ is the relative discriminant, and the
conductor--discriminant formula gives $\mathfrak d_{B_e/B}=\mathfrak f_e$
for the quadratic extension $B_e/B$. Hence
$|\Delta_{B_e}|/|\Delta_B|=241\,\mathrm N\mathfrak f_e=Q_e$. By
Definition~\ref{gf:archimedean}, Proposition~\ref{gf:local-data}(4) and
Lemma~\ref{gf:linear}, if
$\nu_k(e)=0$, the place $v_k$ splits into two real places of $B_e$ and contributes
$\Gamma_\R(s)^2/\Gamma_\R(s)=\Gamma_\R(s)$. If $\nu_k(e)=1$, it becomes one
complex place and contributes $\Gamma_\C(s)/\Gamma_\R(s)=\Gamma_\R(s+1)$.
Thus $\Lambda_{B_e}/\Lambda_B=\Lambda(s,\chi_e)$. It is the completed Hecke
$L$-function of $\chi_e$, which is entire of order at most one
\cite{Hecke1920,Tate1967}.

For (3), the vectors $\bar\iota_1,\bar\iota_2$ are linearly independent. So
$e\mapsto(\nu_1(e),\nu_2(e))$ maps $\Ftwo^8$ onto $\Ftwo^2$ with fibers of
size $64$. Removing $e=0$ from the fiber over $(0,0)$ gives the counts.
\end{proof}

\begin{definition}\label{gf:pure-mixed}
Let $e\ne0$. The \emph{gamma factor} of $\chi_e$ is
$\Gamma_\R(s+\nu_1(e))\Gamma_\R(s+\nu_2(e))$, the factor in the definition of
$\Lambda(s,\chi_e)$ in Proposition~\ref{gf:degree-two}. It is \emph{pure} if
$\nu_1(e)=\nu_2(e)$; writing $\nu$ for the common value, it then equals
$\Gamma_\R(s+\nu)^2$. Otherwise it is \emph{mixed}, and then it equals
$\Gamma_\R(s)\Gamma_\R(s+1)=\Gamma_\C(s)$.
\end{definition}

\begin{corollary}\label{gf:conductors}
For $e\ne0$, $Q_e=241\cdot2^i\cdot m$, where
$2^i=2^{c_{\mathfrak p_1}(e)+c_{\mathfrak p_2}(e)}\in\{1,4,8,16,32,64\}$ and
$m=3^{c_{\mathfrak q_1}(e)+c_{\mathfrak q_2}(e)}5^{c_{\mathfrak r_1}(e)+c_{\mathfrak r_2}(e)}$
divides $225$. The number of characters $\chi_e$ with given $2^i$ and $m$ is
given by the table
\begin{equation}\label{gf:conductor-table}
\begin{array}{c|ccccccccc|c}
 2^i\backslash m&1&3&5&9&15&25&45&75&225&\text{total}\\\hline
 1&0&2&2&1&4&1&2&2&1&15\\
 4&2&4&4&2&8&2&4&4&2&32\\
 8&4&8&8&4&16&4&8&8&4&64\\
 16&1&2&2&1&4&1&2&2&1&16\\
 32&4&8&8&4&16&4&8&8&4&64\\
 64&4&8&8&4&16&4&8&8&4&64\\\hline
 \text{total}&15&32&32&16&64&16&32&32&16&255
\end{array}
\end{equation}
In particular $723\le Q_e\le241\cdot14400=3470400$.
\end{corollary}

\begin{proof}
Only the primes of $S$ divide the conductors
(Proposition~\ref{gf:local-data}(5)), so $Q_e=241\,\mathrm N\mathfrak f_e$
with
\[
 \mathrm N\mathfrak f_e=2^{c_{\mathfrak p_1}(e)+c_{\mathfrak
 p_2}(e)}3^{c_{\mathfrak q_1}(e)+c_{\mathfrak q_2}(e)}5^{c_{\mathfrak
 r_1}(e)+c_{\mathfrak r_2}(e)}.
\]
By Proposition~\ref{gf:local-data},
$c_{\mathfrak p_j}(e)\in\{0,2,3\}$ and
$c_{\mathfrak q_j}(e),c_{\mathfrak r_j}(e)\in\{0,1\}$, and by
Lemma~\ref{gf:linear} these exponents are determined by the vectors of
Table~\ref{tw:vector-table}; counting the $255$ vectors $e\ne0$ gives the
table, which the supplementary program \texttt{check255.gp} also prints
(Remark~\fref{gf:checks}). No nontrivial character has $\mathrm N\mathfrak f_e=1$,
since $B$ has class number one and a unit of norm $-1$
(Lemma~\ref{tw:base}(a),(b)), so that its narrow class number is one. The
range of $Q_e$ can be read off from the table.
\end{proof}

\begin{remark}\label{gf:checks}
The supplementary program \texttt{lfun241.py} (directory
\path{papers/0.04273/certificates} of \cite{NaslundCode4317}) computes, for
each $e\ne0$, the integer $Q_e$, the pair $(\nu_1(e),\nu_2(e))$ and the
coefficients $a_n(e)$ for $n\le N_e=4\lfloor\sqrt{Q_e}\rfloor+1$, from
Propositions~\ref{gf:local-data} and~\ref{gf:degree-two}(1). For the values used here
(Subsection~\ref{an:value-section}), the program \texttt{ye\_hp.py} computes
the coefficients up to the larger truncation point
$8\lfloor\sqrt{Q_e}\rfloor+1$ of that subsection. The supplementary program
\texttt{check255.gp}, in the same directory, compares all $255$ of them with PARI/GP's own number
fields and $L$-functions \cite{PARI}. For each $e$ it computes the maximal
order of $B_e$, checks $|\Delta_{B_e}|=241\,Q_e$ and the pair
$(\nu_1(e),\nu_2(e))$ against the signature of $B_e$, and checks every
coefficient $a_n(e)$ with $n\le N_e$
against the Dirichlet series of $\zeta_{B_e}/\zeta_B$. All agree. It also
prints the table \eqref{gf:conductor-table}, the counts of
Proposition~\ref{gf:degree-two}(3) and the range of $Q_e$. These comparisons
check the implementation; the identification of the $L$-functions rests on
Propositions~\ref{gf:local-data} and~\ref{gf:degree-two}.
\end{remark}
